%! Author = Omar Iskandarani
%! Date = 02/06/2026
%! Affiliation = Independent Researcher, Appingedam, The Netherlands
%! ORCID = 0009-0006-1686-3961
%! DOI = 10.5281/zenodo.xxx

% =========================
% Metadata
% =========================
\newcommand{\papertitle}{A Hierarchical Mass-Energy Functional for Localized Vortex Structures:\\
Kinematic Scale, Geometric Shielding, and Topological State Classification}
\newcommand{\papershorttitle}{Relational Time-of-Arrival}
\newcommand{\paperemail}{info@omariskandarani.com}
\newcommand{\paperlocation}{Appingedam}
\newcommand{\paperkeywords}{relational time, time of arrival, quantum time observables, covariant observables, event currents, clock holonomy, quantum field theory}
\newcommand{\paperdoi}{10.5281/zenodo.xxx}
\newcommand{\paperorcid}{0009-0006-1686-3961}





% ============================================================
%  A Hierarchical Mass-Energy Functional for Localized
%  Vortex Structures: Kinematic Scale, Geometric Shielding,
%  and Topological State Classification
% ============================================================
%\documentclass[aps,prd,preprint,longbibliography,nofootinbib]{revtex4-2}
\documentclass[11pt, a4paper]{article}



% --- STANDARD MATH PACKAGES ---
\usepackage{amsmath,amssymb,amsfonts,bm}
\usepackage{siunitx}
\usepackage[a4paper,margin=1in]{geometry}
\usepackage{silence}
\WarningFilter{latex}{Command \showhyphens has changed}

% =========================
% Core packages (lean)
% =========================
\usepackage[T1]{fontenc}
\usepackage{lmodern}
\usepackage{microtype}
\usepackage{graphicx}
\usepackage{booktabs}
\usepackage{amsthm}
\usepackage{mathtools}
\usepackage{bbm}
\usepackage{enumitem}
\usepackage[colorlinks=true,
            linkcolor=blue!60!black,
            citecolor=blue!60!black,
            urlcolor=blue!60!black]{hyperref}
\usepackage{cleveref}

% ── Theorem environments ─────────────────────────────────────
\newtheorem{definition}{Definition}
\newtheorem{lemma}{Lemma}
\newtheorem{proposition}{Proposition}
\newtheorem{remark}{Remark}

% ── Handy shorthand ──────────────────────────────────────────
\newcommand{\vv}{v}          % characteristic circulation speed
\newcommand{\rc}{r_c}        % core radius
\newcommand{\rhof}{\rho_f}   % background fluid density
\newcommand{\rhoE}{\rho_E}   % kinetic energy density
\newcommand{\rhom}{\rho_m}   % mass-equivalent density
\newcommand{\lc}{\lambda_c}  % Compton wavelength
\newcommand{\Stv}{S_t}       % local time-dilation factor
% ─────────────────────────────────────────────────────────────
\usepackage[absolute,overlay]{textpos}

\usepackage[utf8]{inputenc}
\usepackage{pgfplots}
\pgfplotsset{compat=1.18}
\usetikzlibrary{pgfplots.groupplots}

\begin{document}
    \thispagestyle{empty}%
    \begin{center}
        \Large \bfseries \papertitle \par \vspace{1cm}
        {\Large \itshape \textbf{Omar Iskandarani}\textsuperscript{\textbf{*}} \par} \vspace{0.5cm}
        {\small Independent Researcher, Groningen, The Netherlands\par} \vspace{0.5cm}
        {\today \par}  \vspace{0.5cm}
    \end{center}

    \begin{abstract}
        \normalsize
        We derive and organise a hierarchical mass-energy functional
        $M_K = \rhom\,\rc^3\,\Pi_K\,\Xi_K\,\Stv^{-2}$
        for a localised vortex filament in an ideal incompressible
        fluid, where $\rhom = \frac{1}{2}\rhof v^2/c^2$ is the
        kinetic mass-equivalent density, $\rc$ is the vortex core
        radius, $\Pi_K$ is a dimensionless geometric shielding factor,
        $\Xi_K$ a topological state kernel, and $\Stv$ a local
        time-dilation factor.
        Four results are established without circular reasoning.
        First, the kinetic energy density $\rhoE = \frac{1}{2}\rhof v^2$
        and its mass-equivalent $\rhom = \rhoE/c^2$ follow from
        standard fluid mechanics and special relativity.
        Second, for a vortex filament whose core radius satisfies
        $\rc = \alpha\hbar/(2m_e c)$ (the Compton-scale core hypothesis),
        the dimensionless length ratio $\lc/(\pi\rc) = 4/\alpha$
        follows by pure algebra from the definitions of the Compton
        wavelength $\lc = h/(m_e c)$ and the fine-structure constant
        $\alpha = e^2/(4\pi\varepsilon_0\hbar c)$; no additional
        assumptions are required.
        Third, the local time-dilation factor $\Stv = (1-v^2/c^2)^{1/2}$
        is the standard special-relativistic expression applied to the
        circulation speed.
        Fourth, the topological kernel $\Xi_K$ is explicitly identified
        as an open modelling ansatz whose coefficient values are not
        yet derived from first principles; the paper proposes a
        concrete parameterisation and lists the measurable consequences
        that would confirm or refute it.
        The functional is proposed as an organising framework for
        further analytical and numerical work.
    \end{abstract}
    \vspace{0.5cm}
    \textbf{Keywords:} \paperkeywords


    \begin{textblock*}{1\textwidth}(0.15\textwidth,1.1\textheight)\raggedright\footnotesize\renewcommand{\arraystretch}{1.0} \noindent\rule{\textwidth}{0.4pt}\\[0.5em]
        \textsuperscript{\textbf{*}} Email: \texttt{info@omariskandarani.com}\\
        ORCID: \texttt{\href{https://orcid.org/0009-0006-1686-3961}{0009-0006-1686-3961}}\\
        DOI: \href{https://doi.org/\paperdoi}{\paperdoi}
    \end{textblock*}

    \vfill


% ─────────────────────────────────────────────────────────────


% =============================================================
\section{Introduction and Scope}
\label{sec:intro}
% =============================================================

The mass-energy of a localised fluid structure is determined by
its internal kinetic energy, the spatial extent of the
structure, any geometric factors that modulate the coupling to
the surrounding medium, and the identity of the topological
state.
These four contributions have been studied separately in the
literature on vortex dynamics~\cite{Batchelor1967,Saffman1992},
topological field theory~\cite{FaddeevNiemi1997}, and analogue
gravity~\cite{Barcelo2005}, but have not been assembled into a
single hierarchical formula with explicit status labels for
each factor.

The purpose of this paper is to do exactly that.
We derive each factor in turn, assign it an unambiguous
epistemic label (derived, calibrated, or conjectural), and
assemble the factors into a single mass-energy functional.
The result is an organising framework rather than a predictive
theory: it is predictive only when the topological kernel
$\Xi_K$ is specified by an independent theory of the vortex
topology.

We work throughout in SI units.
Greek indices $\mu,\nu$ run over spacetime; Latin indices
$i,j,k$ over spatial coordinates.
The speed of light in vacuo is $c$; $\hbar = h/(2\pi)$;
$m_e$ and $e$ are the electron mass and charge;
$\varepsilon_0$ is the permittivity of free space.

% =============================================================
\section{Kinematic Scale and Mass-Equivalent Density}
\label{sec:kinematic}
% =============================================================

\subsection{Kinetic energy density of a circulatory flow}
\label{sec:kinetic}

Consider an ideal incompressible fluid of mass density
$\rhof$ (\si{kg.m^{-3}}) supporting a localised circulatory
velocity field $\mathbf{v}(\mathbf{x})$.
The local kinetic energy density is~\cite{Batchelor1967,Lamb1932}
\begin{equation}
  \rhoE(\mathbf{x}) = \tfrac{1}{2}\rhof\,|\mathbf{v}(\mathbf{x})|^2.
  \label{eq:rhoE}
\end{equation}
For a vortex whose peak circulation speed is $\vv$,
the characteristic energy density is
$\rhoE = \frac{1}{2}\rhof\vv^2$.

\subsection{Mass-equivalent density via $E=mc^2$}
\label{sec:rhom}

By the special-relativistic energy--mass equivalence~\cite{Einstein1905},
any energy density $\rhoE$ contributes a mass-equivalent
density
\begin{equation}
  \rhom = \frac{\rhoE}{c^2} = \frac{\rhof\vv^2}{2c^2}.
  \label{eq:rhom}
\end{equation}
This expression is dimensionally consistent:
$[\rhof\vv^2/c^2] = \si{kg.m^{-3}}$.
No assumptions beyond standard fluid mechanics and special
relativity are used.

\subsection{Bare mass scale}
\label{sec:Mstar}

The mass enclosed in a spherical core of radius $\rc$ is
\begin{equation}
  M_\star = \rhom\,\rc^3 = \frac{\rhof\vv^2}{2c^2}\,\rc^3.
  \label{eq:Mstar}
\end{equation}
Here $\rc$ is a free parameter (the vortex core radius) until
it is specified by the topological model; see
\cref{sec:geometric} and \cref{sec:calib}.

% =============================================================
\section{Geometric Shielding Gate}
\label{sec:geometric}
% =============================================================

\subsection{Compton-scale core hypothesis}
\label{sec:core_hypothesis}

\begin{definition}[Compton-scale core]
\label{def:core}
We adopt as a modelling hypothesis the identification
\begin{equation}
  \rc = \frac{\alpha\hbar}{2m_e c},
  \label{eq:rc}
\end{equation}
where $\alpha = e^2/(4\pi\varepsilon_0\hbar c)$ is the
fine-structure constant.
This is a \emph{calibration}, not a derivation: no microscopic
mechanism is asserted here, and the hypothesis is testable by
any experiment that independently constrains the vortex core
radius.
\end{definition}

\subsection{Algebraic derivation of the shielding ratio}
\label{sec:4alpha}

\begin{lemma}[Shielding ratio]
\label{lem:shield}
Under \cref{def:core}, the ratio of the electron Compton
wavelength $\lc = h/(m_e c)$ to the quantity $\pi\rc$ satisfies
\begin{equation}
  \frac{\lc}{\pi\rc} = \frac{4}{\alpha}.
  \label{eq:shielding}
\end{equation}
\end{lemma}

\begin{proof}
The Compton wavelength is $\lc = h/(m_e c) = 2\pi\hbar/(m_e c)$.
Substituting \cref{eq:rc}:
\begin{equation}
  \frac{\lc}{\pi\rc}
  = \frac{2\pi\hbar/(m_e c)}{\pi \cdot \alpha\hbar/(2m_e c)}
  = \frac{2\pi\hbar}{m_e c}\cdot\frac{2m_e c}{\pi\alpha\hbar}
  = \frac{4}{\alpha}.
\end{equation}
The result is an algebraic identity among $\lc$, $\rc$, and
$\alpha$.
It does not depend on the numerical value of any of these
quantities, and introduces no new assumptions beyond
\cref{def:core}.
\end{proof}

\begin{remark}
The ratio $4/\alpha \approx 548.1$ is numerically large because
$\alpha \approx 1/137.036$.
This amplification factor is therefore a direct consequence of
the smallness of the fine-structure constant and the assumption
that the core radius is set at the Compton scale.
\end{remark}

\subsection{Geometric shielding factor}
\label{sec:Pik}

We introduce a binary exposure index $G(K) \in \{0,1\}$ for a
topological state $K$ (see \cref{sec:topo}) and define the
geometric shielding factor
\begin{equation}
  \Pi_K = \left(\frac{\lc}{\pi\rc}\right)^{G(K)}
         = \left(\frac{4}{\alpha}\right)^{G(K)}.
  \label{eq:Pik}
\end{equation}
When $G(K)=0$ the factor is unity (no geometric amplification);
when $G(K)=1$ the factor is $4/\alpha \approx 548$.
The physical interpretation is that states with $G(K)=1$
couple to the medium at the Compton scale, while states with
$G(K)=0$ remain in the bare core regime.
The assignment of $G(K)$ to specific topological states is a
\emph{modelling choice} that lies outside the scope of this
paper.

% =============================================================
\section{Local Time-Dilation Factor}
\label{sec:clock}
% =============================================================

A fluid element moving with speed $v = |\mathbf{v}|$ relative
to the laboratory frame experiences a standard special-relativistic
time dilation~\cite{Will2014,Misner1973}:
\begin{equation}
  \Stv = \sqrt{1 - \frac{v^2}{c^2}},
  \label{eq:St}
\end{equation}
with $d\tau = \Stv\,dt$.
For $v \ll c$ one has $\Stv \approx 1 - v^2/(2c^2)$ and
$\Stv^{-2} \approx 1 + v^2/c^2$.
This factor appears in the mass-energy functional through
the local effective inertia: a region of strong circulation
is both kinetically energetic and locally time-dilated; the
combination amplifies the effective mass.

% =============================================================
\section{Topological State Kernel}
\label{sec:topo}
% =============================================================

\subsection{Motivation and status}
\label{sec:topo_status}

The topological identity of a vortex filament is encoded in
its knot type $K$.
In the mathematical theory of knotted vortices, three families
of invariants are relevant~\cite{Moffatt1969,Ricca2008,Kauffman2001}:
\begin{enumerate}[label=(\roman*)]
  \item \emph{Crossing invariants} $C(K)$, such as the
        crossing number or writhe, which count the complexity of
        the knot diagram;
  \item \emph{Geometric invariants} $L(K)$, such as the
        ropelength~\cite{CantarellaKusner2002}, which encode
        the minimum length required to realise the knot at a
        given core thickness;
  \item \emph{Helicity invariants} $\mathcal{H}(K)$, related
        to the Gauss linking integral and the Hopf
        invariant~\cite{Moffatt1969,Arnold1966}, which measure
        the self-linking of the vorticity field.
\end{enumerate}
A natural dimensionless state kernel is a linear combination
of these invariants.

\subsection{Proposed kernel ansatz}
\label{sec:topo_ansatz}

\begin{definition}[Topological kernel]
\label{def:kernel}
We propose the ansatz
\begin{equation}
  \Xi_K
  = \bigl[\alpha_C\,C(K) + \beta_L\,L(K) + \gamma_H\,\mathcal{H}(K)\bigr]
    \varphi^{-2k(K)},
  \label{eq:Xik}
\end{equation}
where $\alpha_C$, $\beta_L$, $\gamma_H$ are dimensionless
coefficients, $k(K) \in \mathbb{Z}_{\geq 0}$ is a discrete
state-level index, and $\varphi = (1+\sqrt{5})/2$ is the
golden ratio.
\end{definition}

\begin{remark}[Epistemic status of \cref{def:kernel}]
The kernel \eqref{eq:Xik} is a \emph{conjectural modelling
ansatz}, not a derivation.
Three elements are not justified from first principles in this
paper:
\begin{enumerate}[label=(\roman*)]
  \item The coefficients $\alpha_C$, $\beta_L$, $\gamma_H$
        are free parameters; their values must be determined
        by fitting to observed particle masses or by a
        separate dynamical derivation.
  \item The golden-ratio ladder $\varphi^{-2k}$ is motivated
        by the fixed-point structure of the nonlinear
        Schr\"{o}dinger equation on a closed loop
        (cf.~\cite{Pitaevskii1961,Pethick2008}) but is not
        derived from a variational principle for the present
        system.
  \item The binary exposure index $G(K)$ is assigned to knot
        types by a rule not stated in this paper.
\end{enumerate}
These three open points define the research programme that
this paper proposes.
\end{remark}

% =============================================================
\section{The Hierarchical Mass-Energy Functional}
\label{sec:master}
% =============================================================

\subsection{Assembly}
\label{sec:assembly}

Combining \cref{eq:Mstar,eq:Pik,eq:Xik,eq:St}, we define

\begin{proposition}[Hierarchical mass-energy functional]
\label{prop:master}
Under the kinematic, Compton-core, time-dilation, and
topological-kernel hypotheses of
\cref{sec:kinematic,sec:geometric,sec:clock,sec:topo},
the mass-energy of a topological vortex state $K$ at
position $\mathbf{x}$ is
\begin{equation}
  \boxed{
    M_K(\mathbf{x})
    = \frac{\rhof(\mathbf{x})\,|\mathbf{v}(\mathbf{x})|^2}{2c^2}
      \;\rc^3
      \left(\frac{4}{\alpha}\right)^{G(K)}
      \Xi_K
      \left(1-\frac{|\mathbf{v}(\mathbf{x})|^2}{c^2}\right)^{-1}
  }
  \label{eq:master}
\end{equation}
with $\Xi_K$ as in \cref{def:kernel}.
\end{proposition}

\begin{proof}
Substitute \cref{eq:rhom} into \cref{eq:Mstar}, apply
\cref{eq:Pik}, include the state kernel \cref{def:kernel},
and note that $\Stv^{-2} = (1-v^2/c^2)^{-1}$.
The result follows by direct substitution.
The dimensional check gives
$[\rhof v^2/c^2\cdot r_c^3] = \si{kg.m^{-3}.m^3} = \si{kg}$;
all remaining factors are dimensionless.
\end{proof}

\subsection{Factored interpretation}
\label{sec:factors}

The functional separates into four physically distinct sectors:
\begin{equation}
  M_K
  = \underbrace{\vphantom{\Big|}\rhom\,\rc^3}_{\text{medium scale}}
    \times
    \underbrace{\vphantom{\Big|}(4/\alpha)^{G(K)}}_{\text{geometric gate}}
    \times
    \underbrace{\vphantom{\Big|}\Xi_K}_{\text{topological kernel}}
    \times
    \underbrace{\vphantom{\Big|}\Stv^{-2}}_{\text{clock factor}}.
  \label{eq:factored}
\end{equation}
Each factor is independently testable in principle:
the medium scale is set by macroscopic fluid measurements;
the geometric gate follows from the core-radius hypothesis;
the topological kernel depends only on knot invariants;
the clock factor is a standard relativistic correction.

% =============================================================
\section{Reduction Limits}
\label{sec:limits}
% =============================================================

\subsection{Static limit ($v \ll c$)}

For $v \ll c$, $\Stv^{-2} \to 1$ and \cref{eq:master} reduces to
\begin{equation}
  M_K \approx \rhom\,\rc^3\,(4/\alpha)^{G(K)}\,\Xi_K,
\end{equation}
the static mass spectrum.

\subsection{Homogeneous medium}

If $\rhof$ and $v$ are spatially uniform, $\rhom$ is constant
and the spatial integration over the core volume gives a
factor $\propto \rc^3$, consistent with \cref{eq:Mstar}.

\subsection{Weak-circulation expansion}

For $v^2/c^2 \ll 1$,
\begin{equation}
  \Stv^{-2} = 1 + \frac{v^2}{c^2} + \mathcal{O}\!\left(\frac{v^4}{c^4}\right),
\end{equation}
so the leading relativistic correction to the mass is positive
and quadratic in the circulation speed.

% =============================================================
\section{Calibration and Numerical Check}
\label{sec:calib}
% =============================================================

\subsection{Calibrated parameters}
\label{sec:calib_params}

The functional \cref{eq:master} contains three parameters
that are not derived in this paper and must be supplied
externally:
\begin{enumerate}[label=(\roman*)]
  \item $\rhof$: the background effective fluid density.
        We adopt the phenomenological value
        $\rhof = 7.0\times10^{-7}$~\si{kg.m^{-3}},
        consistent with calibrations of the electron vortex
        model against the electron mass~\cite{Iskandarani2026a}.
        This value is a \emph{calibration}, not a prediction.
  \item $\vv$: the characteristic circulation speed.
        We adopt $\vv = \alpha c/2 \approx 1.094\times10^6$~\si{m.s^{-1}},
        the value that identifies the circulation speed with
        $\alpha c/2$~\cite{Iskandarani2026a}.
        Again, this is a \emph{calibration}.
  \item $\rc$: by \cref{def:core},
        $\rc = \alpha\hbar/(2m_e c) \approx 1.409\times10^{-15}$~\si{m}.
\end{enumerate}
The topological coefficients $\alpha_C$, $\beta_L$, $\gamma_H$
and the level index $k(K)$ are open parameters to be
determined.

\subsection{Numerical benchmark}
\label{sec:numerical}

With the calibrated values above:
\begin{align}
  \rhoE &= \tfrac{1}{2}\rhof\vv^2
          \approx 4.188\times10^{5}~\si{J.m^{-3}}, \\
  \rhom &= \rhoE/c^2
          \approx 4.659\times10^{-12}~\si{kg.m^{-3}}, \\
  M_\star &= \rhom\,\rc^3
           \approx 1.303\times10^{-56}~\si{kg}.
\end{align}
The geometric gate evaluates to
\begin{equation}
  \frac{4}{\alpha} \approx 548.14,
\end{equation}
in exact numerical agreement with $\lc/(\pi\rc)$, which
confirms \cref{lem:shield} numerically.
The local time-dilation factor at the calibrated speed is
\begin{equation}
  \Stv = \sqrt{1 - \vv^2/c^2} \approx 0.999993,
  \qquad
  \Stv^{-2} \approx 1.000013,
\end{equation}
so the clock correction is at the level of $1.3\times10^{-5}$
and is negligible in the present calibration regime.

The bare scale $M_\star \approx 1.3\times10^{-56}$~\si{kg}
is much smaller than the electron mass
$m_e \approx 9.1\times10^{-31}$~\si{kg}.
The ratio $m_e/M_\star \approx 7\times10^{25}$ must be supplied
by the product $\Pi_K \Xi_K$, which is the open problem that
the topological kernel \cref{def:kernel} is designed to address.

% =============================================================
\section{Falsifiable Consequences}
\label{sec:falsifiers}
% =============================================================

The functional \cref{eq:master} yields three testable predictions
that are independent of the open parameters
$\alpha_C, \beta_L, \gamma_H$:

\begin{enumerate}[label=(\roman*)]

  \item \textbf{Scaling with core radius.}
        Equation~\eqref{eq:Mstar} predicts $M_\star \propto \rc^3$
        for fixed fluid parameters.
        An independent measurement of $\rc$ that is inconsistent
        with $\rc = \alpha\hbar/(2m_e c)$ would falsify
        \cref{def:core}.

  \item \textbf{Binary geometric clustering.}
        States with $G(K)=1$ should have masses larger than
        states with $G(K)=0$ by a factor of order $4/\alpha \approx 548$,
        holding all other factors equal.
        This is a sharp prediction that does not require knowledge
        of $\Xi_K$.

  \item \textbf{Mass--clock correlation.}
        If the same circulation field $\mathbf{v}(\mathbf{x})$ is
        responsible for both the inertial mass and the time dilation,
        regions of higher effective mass should exhibit a proportionally
        larger proper-time retardation.
        The coupling is $\delta M / M \approx v^2/c^2$ and
        $\delta\tau/\tau \approx -v^2/(2c^2)$ at leading order, giving
        \begin{equation}
          \frac{\delta M}{M} \approx -2\,\frac{\delta\tau}{\tau}.
        \end{equation}
        This relation is independent of all calibrated parameters.

\end{enumerate}

% =============================================================
\section{Status Summary}
\label{sec:status}
% =============================================================

\begin{center}
\begin{tabular}{lll}
\toprule
Factor & Status & Dependencies \\
\midrule
$\rhoE = \tfrac{1}{2}\rhof v^2$ & Derived (fluid mechanics) &
  Batchelor~\cite{Batchelor1967}, Lamb~\cite{Lamb1932} \\
$\rhom = \rhoE/c^2$ & Derived ($E=mc^2$) &
  Einstein~\cite{Einstein1905} \\
$\lc/(\pi\rc) = 4/\alpha$ & Derived (algebra, given \cref{def:core}) &
  \cref{lem:shield} \\
$\Stv = (1-v^2/c^2)^{1/2}$ & Derived (special relativity) &
  Will~\cite{Will2014} \\
$\rc = \alpha\hbar/(2m_e c)$ & Calibration hypothesis & \cref{def:core} \\
$\rhof = 7.0\times10^{-7}$~\si{kg.m^{-3}} & Calibrated &
  \cref{sec:calib_params} \\
$\vv = \alpha c/2$ & Calibrated & \cref{sec:calib_params} \\
$G(K) \in \{0,1\}$ & Modelling rule, not derived & Open \\
$\alpha_C,\beta_L,\gamma_H$ & Free parameters & Open \\
$k(K)$ (golden-ladder level) & Conjectural ansatz & Open \\
\bottomrule
\end{tabular}
\end{center}

% =============================================================
\section{Discussion and Outlook}
\label{sec:discussion}
% =============================================================

The functional \cref{eq:master} separates cleanly into
established physics (kinetic energy density, mass--energy
equivalence, special-relativistic time dilation, algebraic
properties of the fine-structure constant) and open modelling
choices (core-radius hypothesis, topological kernel ansatz,
golden-ladder level assignment).
The boundary between the two is made explicit in
\cref{sec:status}.

The main open problem is the determination of $\Xi_K$ from
a microscopic or variational principle.
Partial progress has been made using the knotted-vortex
energy functionals of Faddeev and Niemi~\cite{FaddeevNiemi1997}
and the ropelength bounds of Cantarella et al.~\cite{CantarellaKusner2002};
a complete derivation remains outstanding.

The core-radius hypothesis \cref{def:core} could in principle
be tested via high-precision measurements of the electron charge
radius or by atomic spectroscopy experiments sensitive to the
inner structure of the Coulomb potential at the femtometre scale.
Existing Lamb-shift data constrain any additional short-range
contribution to the potential at a level that is broadly
consistent with, but does not uniquely select, the value
$\rc = \alpha\hbar/(2m_e c)$~\cite{Bethe1947,Karshenboim2005}.

This paper employs the functional as a toy model for the
purpose of identifying which parts of a localized-vortex mass
formula are generic (following from fluid mechanics and
special relativity) and which parts require a specific
microscopic model.
The three falsifiable predictions of \cref{sec:falsifiers}
hold for any model that shares the kinematic and geometric
sectors, irrespective of the topological kernel.

% =============================================================
\section*{Acknowledgements}
% =============================================================

The author thanks the independent research community for
discussions that helped sharpen the separation between derived
and conjectural elements of this framework.

% =============================================================
\begin{thebibliography}{99}

\bibitem{Batchelor1967}
G.~K. Batchelor,
\textit{An Introduction to Fluid Dynamics},
Cambridge University Press, 1967.
\doi{10.1017/CBO9780511800955}

\bibitem{Saffman1992}
P.~G. Saffman,
\textit{Vortex Dynamics},
Cambridge University Press, 1992.
\doi{10.1017/CBO9780511624063}

\bibitem{Lamb1932}
H.~Lamb,
\textit{Hydrodynamics}, 6th ed.,
Cambridge University Press, 1932.

\bibitem{Einstein1905}
A.~Einstein,
Ist die Tr\"{a}gheit eines K\"{o}rpers von seinem Energieinhalt abh\"{a}ngig?,
\textit{Annalen der Physik} \textbf{18}, 639 (1905).
\doi{10.1002/andp.19053231314}

\bibitem{Moffatt1969}
H.~K. Moffatt,
The degree of knottedness of tangled vortex lines,
\textit{Journal of Fluid Mechanics} \textbf{35}, 117 (1969).
\doi{10.1017/S0022112069000991}

\bibitem{Arnold1966}
V.~I. Arnold,
Sur la g\'{e}om\'{e}trie diff\'{e}rentielle des groupes de Lie de dimension infinie
et ses applications \`{a} l'hydrodynamique des fluides parfaits,
\textit{Annales de l'Institut Fourier} \textbf{16}, 319 (1966).
\doi{10.5802/aif.233}

\bibitem{Ricca2008}
R.~L. Ricca,
Topology bounds energy of knots and links,
\textit{Proceedings of the Royal Society A} \textbf{464}, 293 (2008).
\doi{10.1098/rspa.2007.0174}

\bibitem{FaddeevNiemi1997}
L.~D. Faddeev and A.~J. Niemi,
Stable knot-like structures in classical field theory,
\textit{Nature} \textbf{387}, 58 (1997).
\doi{10.1038/387058a0}

\bibitem{Kauffman2001}
L.~H. Kauffman,
\textit{Knots and Physics}, 3rd ed.,
World Scientific, 2001.
\doi{10.1142/4256}

\bibitem{CantarellaKusner2002}
J.~Cantarella, R.~B. Kusner, and J.~M. Sullivan,
On the minimum ropelength of knots and links,
\textit{Inventiones Mathematicae} \textbf{150}, 257 (2002).
\doi{10.1007/s00222-002-0234-y}

\bibitem{Barcelo2005}
C.~Barceló, S.~Liberati, and M.~Visser,
Analogue gravity,
\textit{Living Reviews in Relativity} \textbf{8}, 12 (2005).
\doi{10.12942/lrr-2005-12}

\bibitem{Will2014}
C.~M. Will,
The confrontation between general relativity and experiment,
\textit{Living Reviews in Relativity} \textbf{17}, 4 (2014).
\doi{10.12942/lrr-2014-4}

\bibitem{Misner1973}
C.~W. Misner, K.~S. Thorne, and J.~A. Wheeler,
\textit{Gravitation},
W.~H. Freeman, 1973.

\bibitem{Pitaevskii1961}
L.~P. Pitaevskii,
Vortex lines in an imperfect Bose gas,
\textit{Soviet Physics JETP} \textbf{13}, 451 (1961).

\bibitem{Pethick2008}
C.~J. Pethick and H.~Smith,
\textit{Bose--Einstein Condensation in Dilute Gases}, 2nd ed.,
Cambridge University Press, 2008.
\doi{10.1017/CBO9780511802850}

\bibitem{Bethe1947}
H.~A. Bethe,
The electromagnetic shift of energy levels,
\textit{Physical Review} \textbf{72}, 339 (1947).
\doi{10.1103/PhysRev.72.339}

\bibitem{Karshenboim2005}
S.~G. Karshenboim,
Precision physics of simple atoms: QED tests, nuclear structure and
fundamental constants,
\textit{Physics Reports} \textbf{422}, 1 (2005).
\doi{10.1016/j.physrep.2005.08.008}

\bibitem{Iskandarani2026a}
O.~Iskandarani,
Resolving three open problems in the swirl-vortex model via the
electron Compton frequency,
Zenodo preprint, 2026.
\doi{10.5281/zenodo.18390075}

\end{thebibliography}

\end{document}